\section*{执行摘要}
\addcontentsline{toc}{section}{执行摘要}

\begin{insight}
\textbf{我们不解释爆款,我们提前发现爆款。}\\[0.3em]
传统社媒工具看「已经火的热搜」,Ripple 看「资本和搜索数据中刚开始升温的早期信号」。借鉴金融界
Digital Oracle 的「资本永远先于舆论」理念,我们用 7 个真实数据源 \textbf{并行扫描} +
CUSUM 算法 + 多 Agent 矛盾推理,在话题登顶热搜的 \textbf{7-14 天前} 就把它交到 KOC 手中。
\end{insight}

\subsection*{产品速览}

\begin{tabularx}{\textwidth}{>{\bfseries\color{ripple-primary}}l X}
\toprule
产品名称 & Ripple(涟漪)\\
\midrule
一句话定位 & KOC 的 Bloomberg Terminal — 在热搜前发现下一个爆款话题 \\
\midrule
核心痛点 & KOC 追热点永远慢半拍,等看到热搜时已是红海(36氪/克劳锐研究) \\
\midrule
核心创新 & 全网首个把「金融早期信号思想」应用到 KOC 内容预测 \\
\midrule
技术架构 & 100\% 借鉴 Claude Code 51 万行代码架构 \\
\midrule
核心模块 & 早期信号雷达(Oracle Agent)+ 12 Agent 内容工厂 \\
\midrule
平台覆盖 & 视频号 / 公众号 / 小红书 / 抖音 / B站 / 微博 全平台支持 \\
\midrule
端覆盖 & Web + macOS/Win 桌面 + 微信小程序 + 移动 PWA \\
\midrule
LLM & 小米 MiMo-V2.5-Pro 首选 + MiniMax M2.7 备用 + 腾讯混元兜底 + BYOK \\
\midrule
合规 & 《AI 生成合成内容标识办法》全合规 + STRIDE 威胁建模 \\
\bottomrule
\end{tabularx}

\vspace{1em}

\subsection*{五大数据亮点(对应大赛评分维度)}

\begin{enumerate}[leftmargin=*]
\item \textbf{创新性} — 7 真实数据源 + CUSUM 算法 + 多 Agent 矛盾推理,业界首创「内容预测的金融化分析框架」
\item \textbf{完整性} — 从「早期信号扫描」到「多平台可发布内容包」全链路闭环,12 个 Agent 协作完成
\item \textbf{用户洞察} — 引用 36 氪「郑州帮」深度报道、CSDN「降 AI 改文言文被喷装」原话等真实痛点
\item \textbf{AI 原生} — 离开多源信号融合 + Agent 推理产品不存在,采用 Claude Code 五层记忆 + 四层压缩纪律
\item \textbf{赛道适配} — 深度集成腾讯混元 API + 元器 Agent + 视频号助手 + 公众号开放平台
\end{enumerate}

\vspace{1em}

\subsection*{与现有竞品的根本差异}

\begin{tabularx}{\textwidth}{l X X}
\toprule
\textbf{维度} & \textbf{现有竞品(千瓜/飞瓜/Buffer/Jasper)} & \textbf{Ripple 涟漪} \\
\midrule
信号视角 & 看已经火的内容 & 看刚开始升温的早期信号 \\
\midrule
预测时效 & T+0 滞后 & T-7 至 T-14 提前发现 \\
\midrule
数据源 & 单一平台抓取 & 7 真实数据源并行 + 矛盾推理 \\
\midrule
AI 形态 & 单 LLM 写文案 & 12 Agent 协作工厂 \\
\midrule
解释能力 & 黑箱推荐 & 每个建议有信号依据 + 置信度 \\
\midrule
平台战场 & 单平台或全平台浅 & 视频号深度 + 全平台覆盖 \\
\bottomrule
\end{tabularx}
